\documentclass[12pt,a4paper]{article}

% Khai báo các gói hỗ trợ tiếng Việt, bảng biểu và toán học
\usepackage[utf8]{inputenc}
\usepackage[T5]{fontenc}
\usepackage[vietnamese]{babel}
\usepackage{geometry}
\usepackage{amsmath}
\usepackage{booktabs}
\usepackage{enumitem}
\usepackage{xcolor}
\usepackage{hyperref}

\geometry{
    a4paper,
    left=25mm,
    right=25mm,
    top=25mm,
    bottom=25mm
}

% Tùy chỉnh màu sắc đường link
\hypersetup{
    colorlinks=true,
    linkcolor=blue,
    filecolor=magenta,      
    urlcolor=cyan,
}

\title{\textbf{BIÊN BẢN NGHIỆM THU: GIAI ĐOẠN 7 -- SINH CÁC CẤU HÌNH HISTORY}}
\author{Dự án: How Much Intelligence Can We Get For How Little?}
\date{\today}

\begin{document}

\maketitle
\hrule
\vspace{15pt}

\section{Mục tiêu}
Giai đoạn 7 tiến hành xây dựng các mức ngân sách lịch sử hành vi $K$ trên tập evaluation cố định. Quá trình này dựa trên cấu hình đã được lựa chọn và đóng băng (freeze) từ giai đoạn phát triển (Giai đoạn 5):
\begin{itemize}[noitemsep]
    \item Cửa sổ thời gian: $W = 1$ ngày;
    \item Giới hạn độ dài lịch sử tối đa: $L_{\max} = 2048$;
    \item Các mức ngân sách thực nghiệm:
    \[
    K \in \{0, 1, 5, 10, 20, 50, 100, \text{Full}\}
    \]
\end{itemize}

\textit{Lưu ý:} Tập test hoàn toàn không được sử dụng trong quá trình xây dựng các cấu hình history này nhằm đảm bảo tính độc lập tuyệt đối.

\section{Phương pháp Xây dựng}
Để tối ưu hóa chi phí tính toán và tránh tiền xử lý dư thừa, thay vì tái dựng history riêng biệt cho từng giá trị $K$, hệ thống tiến hành tái dựng một lần duy nhất lịch sử đầy đủ trong cửa sổ thời gian $W=1$ ngày:
\[
H_t^{(W)} = \{e_i \mid t-W \leq t_i < t\}
\]
Điều kiện $t_i < t$ được áp dụng nghiêm ngặt nhằm đảm bảo không có bất kỳ interaction nào xảy ra tại hoặc sau timestamp của evaluation event bị lọt vào history.

Sau khi xây dựng $H_t^{(W)}$, các cấu hình $K$ hữu hạn được tạo ra bằng phép cắt mảng (slicing), lấy $K$ interaction gần nhất:
\[
H_t^{(K)} = \operatorname{Last}_{\min(K, L_{\max})} \left(H_t^{(W)}\right)
\]

Đối với trường hợp $K=0$ (Video-only baseline), lịch sử là tập rỗng:
\[
H_t^{(0)} = \emptyset
\]

Đối với cấu hình \text{Full}, protocol sử dụng giới hạn $L_{\max}$:
\[
H_t^{(\text{Full})} = \operatorname{Last}_{L_{\max}} \left(H_t^{(W)}\right)
\]
Do đó, cấu hình "Full" không biểu diễn toàn bộ timeline vô hạn của người dùng, mà biểu diễn toàn bộ history hợp lệ trong cửa sổ $W=1$ ngày, sau khi đã áp dụng giới hạn cắt tỉa $L_{\max}=2048$.

\section{Kết quả Xây dựng}
Quá trình sinh lịch sử được thực hiện trên tổng cộng $43,027$ evaluation events. Lịch sử $W=1$ ngày được tái dựng từ $11,713,045$ standard interactions.

\begin{table}[h]
\centering
\caption{Thống kê độ dài các mức ngân sách lịch sử $K$}
\label{tab:phase7-history-stats}
\begin{tabular}{lrrrr}
\toprule
\textbf{$K$} & \textbf{Mean length} & \textbf{Median} & \textbf{Max} & \textbf{Zero history} \\
\midrule
0    & 0.00   & 0   & 0    & 43,027 \\
1    & 1.00   & 1   & 1    & 93 \\
5    & 4.98   & 5   & 5    & 93 \\
10   & 9.95   & 10  & 10   & 93 \\
20   & 19.83  & 20  & 20   & 93 \\
50   & 48.86  & 50  & 50   & 93 \\
100  & 94.82  & 100 & 100  & 93 \\
Full & 574.55 & 419 & 2048 & 93 \\
\bottomrule
\end{tabular}
\end{table}

Phân tích dữ liệu:
\begin{itemize}
    \item Lịch sử $H_t^{(1\text{day})}$ ban đầu (trước khi áp dụng $L_{\max}$) có độ dài trung bình là $603.15$, trung vị $419$ và giá trị lớn nhất (maximum) lên tới $10,259$. 
    \item Sau khi áp dụng giới hạn ngân sách Full ở mức $L_{\max}=2048$, giá trị maximum giảm xuống đúng $2048$ và độ dài trung bình giảm còn $574.55$.
    \item Có đúng $93$ evaluation events không có bất kỳ interaction hợp lệ nào trong cửa sổ 1 ngày. Các sự kiện này \textbf{vẫn được giữ lại} trong evaluation set và được biểu diễn bằng history rỗng, tuân thủ tuyệt đối quy tắc xử lý thiếu history của protocol.
\end{itemize}

\section{Lưu trữ Kết quả (Artifacts)}
Toàn bộ các cấu hình history đã được lưu trữ độc lập tại thư mục:
\begin{verbatim}
data/processed/phase7/
\end{verbatim}
Danh sách các tệp dữ liệu bao gồm:
\begin{itemize}[noitemsep]
    \item \texttt{history\_K0.parquet}
    \item \texttt{history\_K1.parquet}
    \item \texttt{history\_K5.parquet}
    \item \texttt{history\_K10.parquet}
    \item \texttt{history\_K20.parquet}
    \item \texttt{history\_K50.parquet}
    \item \texttt{history\_K100.parquet}
    \item \texttt{history\_KFull.parquet}
    \item \texttt{history\_config.yaml}
    \item \texttt{history\_config\_summary.csv}
\end{itemize}

\section{Kiểm tra Protocol (Audit)}
Các invariant cốt lõi của Giai đoạn 7 đã được xác nhận:
\begin{itemize}[noitemsep]
    \item Cấu hình $W=1$ ngày: \textbf{FROZEN}
    \item Cấu hình $L_{\max}=2048$: \textbf{FROZEN}
    \item Evaluation set cố định giữa các mức $K$: \textbf{PASS}
    \item Không rò rỉ dữ liệu từ tập Test: \textbf{PASS}
    \item Tối ưu hóa tính toán (xây dựng history 1 lần): \textbf{PASS}
    \item Sinh các mức $K$ bằng phương pháp Slicing: \textbf{PASS}
\end{itemize}

\vspace{10pt}
\begin{center}
\large\textbf{KẾT LUẬN: GIAI ĐOẠN 7 -- HOÀN THÀNH 100\%}
\end{center}
\noindent Dữ liệu của tất cả các ngân sách $K$ đã được đóng băng và hoàn toàn sẵn sàng để đưa vào pipeline huấn luyện mô hình (Giai đoạn 8 \& 9).

\end{document}